\documentclass[11pt]{article}
\usepackage[a4paper,margin=28mm]{geometry}
\usepackage[T1]{fontenc}
\usepackage[utf8]{inputenc}
\usepackage{lmodern,microtype}
\usepackage{amsmath,amssymb,amsthm,mathtools}
\usepackage[hidelinks]{hyperref}
\hypersetup{pdftitle={A joint moment bound for the three largest graph eigenvalues},pdfauthor={}}
\usepackage{enumitem}
\newtheorem{theorem}{Theorem}[section]
\newtheorem{lemma}[theorem]{Lemma}
\newtheorem{proposition}[theorem]{Proposition}
\newtheorem{corollary}[theorem]{Corollary}
\theoremstyle{remark}\newtheorem{remark}[theorem]{Remark}
\newcommand{\E}{\mathbb E}
\newcommand{\R}{\mathbb R}
\newcommand{\cstar}{c_{\mathrm{w}}}
\DeclareMathOperator{\tr}{tr}
\DeclareMathOperator{\diag}{diag}
\DeclareMathOperator{\supp}{supp}
\DeclareMathOperator{\rank}{rank}
\DeclareMathOperator{\esssup}{ess\,sup}
\DeclareMathOperator{\essinf}{ess\,inf}
\allowdisplaybreaks[1]
\numberwithin{equation}{section}
\title{A joint moment bound for the three largest graph eigenvalues}
\author{}
\date{}
\begin{document}
\maketitle
\begin{abstract}
Let $\lambda_1(G)\ge\lambda_2(G)\ge\cdots$ be the adjacency eigenvalues
of a finite simple graph. We prove
$\lambda_1(G)+\lambda_2(G)+\lambda_3(G)\le cn-3$ for every graph of order
$n\ge3$, where $c=1.284673255723918505\ldots$ is an explicit algebraic
constant. This improves the coefficient $(3+\sqrt5)/4$ in the universal
bound of Sivashankar, Tang and Wakhare. The proof controls the signed
mean and the absolute Gram energy of a single isotropic law by a joint
third- and fourth-order moment identity. Radial and directional errors
are retained in the same identity. We also prove that the canonical
eigenfunction law of every maximizing graphon has at most 24 support
points. The structural argument uses moment-preserving variations and
a separating family of binary observations to rule out minimal
second-moment dependencies on positive cliques. A small final improvement
comes from positive-mass self-pairs. Neither the optimal coefficient nor
the extremal graphs are determined.
\end{abstract}

\noindent\textbf{Keywords.} Adjacency eigenvalues; spectral sums; graphons;
isotropic measures; projection constants.

\section{Introduction}
For a finite simple undirected graph $G$ of order $n\ge3$, put
\[
 \Lambda_3(G)=\lambda_1(G)+\lambda_2(G)+\lambda_3(G),
\]
where the adjacency eigenvalues are ordered algebraically. The signs in
this sum are retained. The extremal problem asks how large $\Lambda_3(G)$
can be when only $n$ is prescribed. The spectral-sum problem appears in Ebrahimi's thesis~\cite{Ebrahimi};
general bounds for sums of leading eigenvalues were developed by
Mohar~\cite{Mohar}. More recently,
Sivashankar, Tang and Wakhare~\cite[Proposition 6.1]{STW} proved
\begin{equation}\label{stw-bound}
 \Lambda_3(G)\le\frac{3+\sqrt5}{4}n-3.
\end{equation}
Their argument uses the rank-three absolute projection constant
$\gamma(3)=(1+\sqrt5)/2$; a proof of the requisite projection bound is
provided by Der\k{e}gowska and Lewandowska~\cite{DL}.

The two quantities in that reduction are naturally coupled. If
$X\in\R^3$ satisfies $\E XX^T=I_3$ and $X'$ is an independent copy, they are
\[
 A=\E|X\cdot X'|,\qquad B=|\E X|^2.
\]
The separate estimates $A\le(1+\sqrt5)/2$ and $B\le1$ give
\eqref{stw-bound}. A better graph bound therefore requires a quantitative
restriction on their simultaneous size. The mixed projection constant
expressing this question was already introduced in~\cite[Section 6.2]{STW}.
We work directly with its isotropic-law formulation.

Define
\begin{equation}\label{main-constant}
 c=\frac{106106902139519}{107310156800000}
       +\frac{41\sqrt3}{240}-\frac{53}{9000000}
   =1.284673255723918505\ldots.
\end{equation}
\begin{theorem}\label{main}
Every finite simple unweighted graph $G$ on $n\ge3$ vertices satisfies
\begin{equation}\label{main-bound}
 \Lambda_3(G)\le cn-3.
\end{equation}
Moreover, every Borel probability law on $\R^3$ satisfying
$\E XX^T=I_3$ obeys
\begin{equation}\label{law-bound}
 \E|X\cdot X'|+|\E X|^2\le2c.
\end{equation}
Only a finite second moment is assumed; zero vectors and unbounded radii
are allowed.
\end{theorem}
The coefficient decrease from~\eqref{stw-bound} is
$0.024343738651028919\ldots$, with the constant term $-3$ unchanged.
This comparison concerns the stated universal order-only bound. It does
not compare all estimates involving additional graph parameters, nor
assert that $c$ is optimal.

We use graphons to pass to an attained maximum. For a symmetric measurable
kernel $W:[0,1]^2\to[0,1]$, let
\begin{equation}\label{sig-definition}
 \Sigma_k(W)=\sup_{\E FF^T=I_k}
       \iint W(x,y)F(x)\cdot F(y)\,dx\,dy,
 \qquad \mathcal M=\sup_W\Sigma_3(W).
\end{equation}
The supremum is over real orthonormal $k$-frames. Since $T_W$ is compact
and self-adjoint on an infinite-dimensional space, $\Sigma_k$ is the sum
of the first $k$ positive eigenvalues, padded by zero if necessary. This
convention applies only to graphons, not to $\Lambda_3(G)$.
The maximum $\mathcal M$ is attained. For a maximizing graphon, choose an
orthonormal top-three eigenframe $F=(f_1,f_2,f_3)^T$, with $f_1\ge0$,
and call $F_\#dx$ its canonical law.

\begin{theorem}\label{structure-main}
Let $W$ attain $\mathcal M$, let $\sum_i p_i\delta_{v_i}$ ($p_i>0$) be its
canonical law, and let $L$ be the diagonal matrix of the three leading
eigenvalues of $W$. Then:
\begin{enumerate}[leftmargin=2.2em,itemsep=2pt]
\item[(a)] the support has at most $24$ points;
\item[(b)] $W$ agrees almost everywhere with the binary step kernel
$\mathbf1_{\{v_i\cdot v_j>0\}}$ on the corresponding positive-measure
fibers;
\item[(c)] the support points satisfy
\begin{gather*}
 \sum_i p_iv_iv_i^T=I_3,\qquad v_i^TLv_i=\mathcal M,\qquad v_{i1}>0,\\
 v_i\cdot v_j\ne0,\qquad
 B_{ij}=\mathbf1_{\{v_i\cdot v_j>0\}},\qquad
 \sum_jp_jB_{ij}v_j=Lv_i;
\end{gather*}
\item[(d)] for every positive-dot clique $J$ in the support, the matrices
$\{v_iv_i^T:i\in J\}$ are linearly independent; in particular $|J|\le6$.
\end{enumerate}
\end{theorem}
This assertion concerns every original maximizer, not merely the
existence of a maximizer with small support.

\subsection{Method and relation to earlier work}
The proof has two parts. First, second variation on the actual probability
space gives a negative semidefinite kernel. On each of four sign classes,
this controls the rank of the positive-part Gram kernel and forces finite
canonical support. Moment-preserving variations then imply vector
relations against every binary row of the maximizing kernel. Those rows
separate distinct support points. A minimal-dependency argument shows
that no second-moment relation can survive on a positive clique.

The variation, cloning, and orthant-reduction mechanisms have close
antecedents in Kumar, Mohar, Mojallal and Pragada~\cite{KMMP}. Related graphon
ellipse and clique-cloning arguments for the two-eigenvalue problem
appear in Kumar, Liu, Monterde, Pragada and Tait~\cite{KLMPT}. General
bounds for minimal dependencies of powers of linear forms are due to
Li and Liu~\cite{LiLiu}. Our separating-observation lemma is stated and
proved with its specific hypotheses in Section~\ref{sec-separation};
we do not claim a new general bound for such dependencies.

Second, for one isotropic law we retain the mean deficit, radial variance,
regression error, and weighted harmonic moments of degrees three and four.
A signed quartic majorant couples them in one exact identity. Directional
Hermite inequalities and a radial comparison lead to an explicit sum of
nonnegative terms. Harmonic moment and polynomial-majorant methods are
classical; see Glazyrin~\cite{Glazyrin}, Der\k{e}gowska and
Lewandowska~\cite{DL,DL-SW}, and Sivashankar, Tang and Wakhare~\cite{STW}.
The contribution here is the particular joint control of these errors
for the mixed signed and absolute energy, rather than the use of
polynomial majorants by itself. The support theorem supplies a small
additional term through self-pairs of actual atoms. Diagonal removal and
effective-cardinality estimates are standard weighted-energy devices;
see~\cite{Weighted}.

Sections~\ref{sec-variational}--\ref{sec-separation} prove the structural
result. Sections~\ref{sec-moments}--\ref{sec-reserve} give the joint
moment certificate. Section~\ref{sec-conclusion} applies it to a genuine
maximizer and transfers the bound to all isotropic laws and finite graphs.

\section{Variational structure of a maximizing graphon}\label{sec-variational}
We may work on any standard nonatomic probability space. All changes of
measure below can be realized on such a space.
\begin{lemma}\label{ranktwo}
Every $[0,1]$-valued graphon satisfies $\Sigma_2(W)\le7/6$.
\end{lemma}
\begin{proof}
We give the elementary argument for the rank-two absolute-energy
bound needed below.
Take any $\R^2$-valued law $\E YY^T=I_2$, write $Y=RU$ with a unit
choice at zero, $C=\E R$, and $N=\E[R UU^T]$.
Then $C^2\le2$ and
$\tr(N^2)\le1$, the latter from
$\E\|U-NY\|^2=1-\tr(N^2)\ge0$.
On the circle,
$g_{4,2}(t)=t^4-t^2+1/8$ is a positive kernel: if
$U=(\cos a,\sin a)$, it is $\cos(4(a-a'))/8$.
Hence its $RR'$-weighted integral is nonnegative.
The exact majorant
\[
 \frac29+\frac{11}{9}t^2-\frac49t^4-|t|
     =\frac{(1-|t|)(2|t|-1)^2(|t|+2)}9\ge0
\]
therefore gives
\[
 \E|Y\cdot Y'|\le\frac5{18}C^2+\frac79\tr(N^2)\le\frac43.
\]
Also $|\E Y|^2\le1$. For any orthonormal two-frame in a graphon,
$W(Y\cdot Y')\le(Y\cdot Y')_+$ pointwise, so its trace is at most
$(4/3+1)/2=7/6$. Taking the supremum proves the lemma.
The absolute-energy estimate is classical; see~\cite[Theorems 2 and 4]{Glazyrin}.
\end{proof}


\begin{lemma}\label{attainment}
The supremum $\mathcal M$ is attained, and
\[
 \cstar:=\frac{9+\sqrt5+2\sqrt{35}}{18}\le\mathcal M
       \le\frac{1+\sqrt3}{2}.
\]
Every maximizing graphon has three positive leading eigenvalues.
\end{lemma}
\begin{proof}
For an orthonormal $k$-frame, $K(x,y)=F(x)\cdot F(y)$ satisfies
$\int K=|\int F|^2\le1$ and $\int K^2=k$. Hence
$\int WK\le\int K_+\le(1+\sqrt k)/2$.
For the lower bound, use a six-part kernel with template $K_1\vee C_5$
and diagonal blocks equal to one. Give the center mass $a=1-5b$ and each
rim part mass $b$, where $0<b<1/5$. Put $\phi=(1+\sqrt5)/2$.
Its quotient spectrum is
\[
 \frac{1-2b\pm\sqrt{1+4b-36b^2}}2,
 \quad\phi b,\ \phi b,\ -b/\phi,\ -b/\phi.
\]
Indeed the four nonconstant Fourier modes on the pentagon give the last
four values; the remaining block is
$\left(\begin{smallmatrix}a&\sqrt{5ab}\\\sqrt{5ab}&3b\end{smallmatrix}\right)$.
Its positive eigenvalue exceeds $3b$, and its determinant is $-2ab<0$.
Consequently the sum of its three positive eigenvalues is
\[
 \frac{1+2\sqrt5 b+\sqrt{1+4b-36b^2}}2.
\]
Substitution of $b=(7+5\sqrt7)/126$ gives $\cstar$.
This is the usual quotient calculation for a join construction;
see~\cite{Hjoin}. In particular $\cstar>32/25>7/6$.

For completeness we record why $\Sigma_3$ of this step kernel is exactly the
sum of the three positive quotient eigenvalues. If $f$ is supported in one
part $A_i$ and $\int_{A_i}f=0$, then for $x\in A_i$ one has
$(T_Wf)(x)=\int_{A_i}f=0$, while for $x$ in any other part $A_j$ one has
$(T_Wf)(x)=W_{ji}\int_{A_i}f=0$; hence $T_Wf=0$. The whole block-mean-zero
subspace therefore lies in the zero eigenspace, and $\Sigma_3$ is governed
by the quotient matrix alone.

For attainment, let $D=W_1-W_2$. If $\lambda\ne0$ is an eigenvalue of
$T_D$ with unit eigenfunction $f$, then $|D|\le1$ gives
$\|f\|_\infty\le1/|\lambda|$. The cut-norm dual estimate yields
\[
 |\lambda|=|\langle f,T_Df\rangle|
 \le4\|D\|_\square\|f\|_\infty^2,
 \qquad \|T_D\|_{2\to2}^3\le4\|D\|_\square.
\]
The assertion is trivial when $T_D=0$; otherwise compactness supplies an
eigenvalue of modulus $\|T_D\|$. By the variational definition,
$|\Sigma_3(W_1)-\Sigma_3(W_2)|\le3\|T_D\|$. This also holds after
relabeling. Cut-metric compactness~\cite{LS} therefore gives a maximizer.
Finally Lemma~\ref{ranktwo} implies
$\lambda_3\ge\cstar-7/6>0$ at each maximizer.
\end{proof}

Fix a maximizing $W$ and a top-three orthonormal eigenframe $F$, taking
$f_1\ge0$. Write $L=\diag(\lambda_1,\lambda_2,\lambda_3)$.
Taking the absolute value of a top Rayleigh maximizer supplies the
nonnegative choice of $f_1$. We do not initially assume any eigenvalue simplicity.
\begin{lemma}\label{ellipse-structure}
Almost everywhere, $F^TLF=\mathcal M$. Moreover
$\lambda_3>\lambda_4$, where the fourth positive eigenvalue is zero if
absent. The canonical support $S=\supp(F_\#dx)$ is compact, lies on the
ellipsoid $v^TLv=\mathcal M$, and avoids the origin.
\end{lemma}
\begin{proof}
For a bounded real $g$ of mean zero, change density to $1+tg$ for both
signs of sufficiently small $t$. Set $Q=\int gFF^T$. The unchanged frame
has complete Gram matrix $I+tQ$ and numerator matrix
$L+t(QL+LQ)+O(t^2)$. After a common orthonormalization its trace is
\[
 \tr\bigl[(I+tQ)^{-1}\{L+t(QL+LQ)+O(t^2)\}\bigr]
   =\mathcal M+t\tr(LQ)+O(t^2).
\]
Maximality gives $\int gF^TLF=0$. Since every bounded mean-zero $g$ is
allowed, $F^TLF$ equals its mean $\mathcal M$ almost everywhere.
This applies to every active eigenframe. If
$\lambda_3=\lambda_4=\lambda>0$, compare the frames with third entries
$f_3$, $f_4$ and $(f_3+f_4)/\sqrt2$, keeping $f_1,f_2$ fixed.
The three identities imply $f_3^2=f_4^2$ and $f_3f_4=0$ almost everywhere,
a contradiction. Only finitely many null sets are combined.
Since $L\succ0$, the ellipsoid is compact and excludes zero. It contains
the full support, not just almost every frame value.
\end{proof}

\begin{lemma}\label{row-insertion}
For every measurable $h:[0,1]\to[0,1]$, the vector $r=\int hF$ satisfies
$r^TL^{-1}r\le\mathcal M$.
\end{lemma}
\begin{proof}
Give the original space $A$ mass $u=1-\epsilon$ and a new independent part
$B$ mass $\epsilon$; the kernel is $W$ on $A\times A$, equals the prescribed
cross-row $h$ (a function on $A$) on $A\times B$ and $B\times A$, and is zero
on $B\times B$. Put $t=L^{-1}r$ and $T=|t|^2$. For the trial frame take $f_i$
on $A$ and the constant $t_i$ on $B$; then the $A\times B$ and $B\times A$
contributions to the numerator are $\epsilon u\,r_it_j$ and its transpose.
Dividing the entire frame by $\sqrt u$ leaves every Rayleigh quotient
unchanged, and the resulting Gram and numerator matrices are
\[
 G=I+(\epsilon/u)tt^T,\qquad
 N=uL+\epsilon(rt^T+tr^T).
\]
The rank-one inverse formula gives the exact trial trace
\[
 \tr(G^{-1}N)=u\mathcal M+
       \frac{\epsilon r^TL^{-1}r}{1+\epsilon T/u}.
\]
This cannot exceed $\mathcal M$. Divide by $\epsilon$ and let
$\epsilon\downarrow0$.
\end{proof}

\begin{lemma}\label{contact-structure}
For $\nu=F_\#dx$ and every $v\in S$,
\begin{equation}\label{all-rows}
 \nu\{w:v\cdot w=0\}=0,\qquad
 \int_{v\cdot w>0}w\,d\nu(w)=Lv,\qquad v_1>0.
\end{equation}
The original kernel is $\mathbf1_{\{F(x)\cdot F(y)>0\}}$ almost
everywhere. It is irreducible, and $\lambda_1>\lambda_2$.
\end{lemma}
\begin{proof}
Let $Z=\{\int hF:0\le h\le1\}$ and
$E=\{r:r^TL^{-1}r\le\mathcal M\}$. The set $Z$ is compact by weak-star
compactness and is contained in $E$ by Lemma~\ref{row-insertion}.
Its support function is $H(v)=\int(v\cdot F)_+$.
For almost every $x$, the original row has integral vector $LF(x)$,
and $F(x)^TLF(x)=\mathcal M$. Since the support function of $E$ is
$\sqrt{\mathcal M v^TLv}$, it follows that $H(F(x))=\mathcal M$.
Continuity then gives $H(v)=\mathcal M$ for every $v\in S$.

For such $v$, the unique support point of the strictly convex ellipsoid
$E$ in direction $v$ is $Lv$. The row $h_v=\mathbf1_{\{v\cdot F>0\}}$
attains this support value, so $\int h_vF=Lv$. If
$A_0=\{v\cdot F=0\}$, then $h_v+\mathbf1_A$ has the same support value
for every measurable $A\subset A_0$. Uniqueness forces $\int_A F=0$
for every such $A$, hence $F=0$ almost everywhere on $A_0$.
The ellipsoid identity implies that $A_0$ is null.

It remains to identify every maximizing row, not merely those obtained by
perturbing $h_v$ on $A_0$. Let $h\in[0,1]$ satisfy $\int hF=Lv$ for some
$v\in S$. Pairing with $v$ and using $v^TLv=\mathcal M$ gives
$v\cdot\int hF=\mathcal M$, while on the other hand
\[
 v\cdot\int hF=\int h\,(v\cdot F)\,dx\le\int(v\cdot F)_+\,dx=H(v)=\mathcal M .
\]
Hence equality holds throughout. Since $0\le h\le1$, the pointwise bound
$h\,(v\cdot F)\le(v\cdot F)_+$ is therefore an equality almost everywhere,
which forces $h=1$ almost everywhere on $\{v\cdot F>0\}$ and $h=0$ almost
everywhere on $\{v\cdot F<0\}$. As $A_0$ is null, $h=h_v$ almost
everywhere. Thus the maximizing row is unique modulo null sets, which
identifies every original row and proves the threshold formula.

If a nontrivial measurable cut had zero cross-kernel, its indicator
multiplier would commute with $T_W$ and with the isolated top-three
spectral projection. The kernel $F(x)\cdot F(y)$ of that projection
would vanish across the cut, contradicting the zero-equator assertion
and Fubini. Thus $W$ is irreducible. The zero set of the nonnegative
$f_1$ would also give a zero cut by its eigen-equation, so $f_1>0$ almost
everywhere. If another top eigenfunction orthogonal to $f_1$ existed,
it would change sign. Equality in the absolute-value Rayleigh comparison
would force zero edges between its positive and negative sets; its
absolute value is a nonnegative top eigenfunction and is positive almost
everywhere. This contradicts irreducibility, proving simplicity.
Finally, for every $v\in S$, $\lambda_1v_1=\int h_vf_1>0$:
otherwise $h_v=0$ almost everywhere and hence $Lv=0$, which is impossible.
\end{proof}
\section{Finite canonical support}\label{sec-finite}
Continue with the maximizing graphon and canonical law of the previous
section, and write $h_v(w)=\mathbf1_{\{v\cdot w>0\}}$. Define continuous symmetric kernels on $S\times S$ by
\[
 A(v,w)=(v\cdot w)_+,\qquad
 D(v,w)=(v\cdot w)(v^TLw),\qquad K=A-D.
\]
\begin{theorem}[Second variation on the actual support]\label{negative}
For every finite list $v_1,\ldots,v_s\in S$ and every real list $a_1,\ldots,a_s$,
\begin{equation}
 \sum_{i,j}a_i a_j K(v_i,v_j)\le0. \label{nsd}
\end{equation}
The assertion concerns only the actual canonical support.
\end{theorem}
\begin{proof}
For any bounded measurable real $g$ on $S$ with $\int g\,d\nu=0$, use the legitimate density
$1+tg(F(x))$ for both signs of sufficiently small $t$.
Set
\[
 Q=\int g(v)vv^T\,d\nu(v),\qquad
 R=\iint h_v(w)vw^Tg(v)g(w)\,d\nu(v)d\nu(w).
\]
The unchanged three coordinate functions have the complete Gram $G_t=I+tQ$
and numerator matrix $N_t=L+t(QL+LQ)+t^2R$ in this changed measure.
For small $t$, $G_t$ is positive definite; multiplying the frame by $G_t^{-1/2}$
is a genuine common orthonormalization. Since
$N_t=G_tLG_t+t^2(R-QLQ)$ and $\tr(LQ)=\mathcal M\int g=0$, its exact trace is
\begin{equation}
 \tr(G_t^{-1}N_t)=\mathcal M+t^2\tr\bigl(G_t^{-1}(R-QLQ)\bigr)\le \mathcal M. \label{density}
\end{equation}
Letting $t\to0$ gives $\iint K(v,w)g(v)g(w)\,d\nu(v)d\nu(w)\le0$.
There is no discarded remainder or varying-scheme constant in this calculation.
By Lemma~\ref{ellipse-structure} and~\eqref{all-rows}, $\int K(v,w)d\nu(w)=\mathcal M-\mathcal M=0$ for every $v\in S$.
Subtracting the mean therefore proves the same inequality for every bounded $g$.
To obtain \eqref{nsd}, combine repeated points and take disjoint shrinking
neighborhoods of the remaining points. Each has positive $\nu$-mass by full support.
Use $g=\sum_i a_i1_{U_i}/\nu(U_i)$, separately for each neighborhood size.
This auxiliary $g$ is bounded but need not have mean zero. The legal probability
density direction is $g-\int g\,d\nu$; centering leaves the $K$ quadratic form unchanged.
Its bound need not be uniform across neighborhood sizes.
Continuity of $K$ on compact $S\times S$ gives \eqref{nsd} in the limit.
This limiting step needs no uniform permissible size of $t$.
No eigenvalue differentiation is used, including when $\lambda_2=\lambda_3$.
\end{proof}

\subsection{Propagation of moment relations}
\begin{lemma}\label{propagate}
If $v_i\in S$ have $v_i\cdot v_j>0$ for all $i,j$, then
\[
 \sum_i a_i v_iv_i^T=0
 \quad\Longrightarrow\quad
 \sum_i a_i v_i=0\quad\hbox{and}\quad
 \sum_i a_i A(v_i,w)=0\quad\hbox{for every }w\in S.
\]
Consequently the vector space spanned by the functions $A(v,\cdot)$, $v\in S$,
has dimension $r\le24$ and admits a basis $A(z_j,\cdot)$ with actual $z_j\in S$.
\end{lemma}
\begin{proof}
Write $Q_a=\sum a_i v_iv_i^T$ and $b_a=\sum a_i v_i$.
The quadratic form in \eqref{nsd} is $|b_a|^2-\tr(LQ_a^2)$ on this positive-dot set.
For $Q_a=0$ it is nonnegative, hence zero and $b_a=0$.
A null vector of a negative semidefinite quadratic form has zero cross terms:
apply \eqref{nsd} also to this list with $w$ appended and coefficient $s\in\R$.
The resulting polynomial $2s\sum_i a_iK(v_i,w)+s^2K(w,w)$ is nonpositive for all $s$,
so its linear coefficient vanishes. The $D$ sum vanishes because $Q_a=0$;
therefore the $A$ sum vanishes for every $w$.
Partition $S$ by the signs of coordinates 2 and 3, assigning zero to the positive sign.
Within each of the four disjoint sets, every dot product is strictly positive,
since coordinate 1 is positive and the other two coordinate products are nonnegative.
In each nonempty set choose a basis, of size at most 6, of its matrices $vv^T$
in the six-dimensional space of symmetric $3$ by $3$ matrices.
The implication just proved expresses every corresponding $A$ row in the span
of those basis-point rows. Combine the four lists and discard dependent rows.
This proves $r\le24$. These are spanning witnesses, not a replacement vertex measure.
This bounds the row space, not yet the cardinality of the support.
\end{proof}

\subsection{Sign regions and finite support}
\begin{lemma}[Sign-region dimension descent]\label{descent}
Let $T\subset S$ have linear span of dimension $d\ge1$.
Suppose for every $w\in S$ there is a vector $b_w\in\R^3$ such that
$A(v,w)=v\cdot b_w$ for all $v\in T$. Then $|T|\le2^{d-1}$.
\end{lemma}
\begin{proof}
If $d=1$, the positive first coordinate and $v^TLv=\mathcal M$ make $T$ a singleton.
For general $d$, if no $w\in S$ has strictly positive dot product with one member
of $T$ and strictly negative dot product with another, any two members have the
same threshold row outside the union of their two zero equators.
Equation \eqref{all-rows} then makes the members equal.
Otherwise choose such a $w$, with witnesses $v_+,v_-$, and split
$T_+=\{v\in T:v\cdot w\ge0\}$ and $T_-=\{v\in T:v\cdot w<0\}$.
On $T_+$, $v\cdot(b_w-w)=0$; on $T_-$, $v\cdot b_w=0$.
Both are proper hyperplanes of $\operatorname{span}T$: evaluate $b_w-w$ at $v_-$
and $b_w$ at $v_+$. Thus both children have span dimension at most $d-1$.
They inherit the linearization for every $w'\in S$.
Induction gives at most $2^{d-2}$ points in each child, proving the bound.
Zeros were included in $T_+$, so none are lost in the dimension descent.
\end{proof}

\begin{theorem}[Finite support of every maximizer]\label{finite}
For every global maximizer of $\Sigma_3$, its actual canonical support satisfies
\begin{equation}
 |S|\le 4\bigl(1+r(r+1)/2\bigr)\le1204=:N_*. \label{bound}
\end{equation}
The original graphon agrees almost everywhere with a binary step graphon
on these actual fibers, with
positive masses, diagonal blocks 1, and no equal rows. For this particular
maximizer $\min_{v,w\in S}|v\cdot w|>0$.
\end{theorem}
\begin{proof}
Take the actual anchors $z_1,\ldots,z_r$ from Lemma~\ref{propagate}.
Symmetry gives, for each $w\in S$, coefficients $c_j(w)$ with
$A(v,w)=\sum_j c_j(w)A(v,z_j)$ for every $v\in S$.
The union of the finitely many anchor equators has $\nu$-measure zero by \eqref{all-rows}.
On any fixed strict sign region of these anchors,
$A(v,z_j)$ is either $v\cdot z_j$ or zero. Thus $A(v,w)$ is linear in $v$
there, for every $w$, with coefficients determined by that region.
Because $v_1>0$, use coordinates $(v_2/v_1,v_3/v_1)$.
Each anchor equator becomes an affine line, or is empty if its last two coordinates vanish.
At most $r$ lines have at most $1+r(r+1)/2$ complementary regions: the $j$-th
line has at most $j-1$ distinct previous intersection points and adds at most $j$ regions.
Coincident or parallel lines only lower this count. A strict sign pattern is an
intersection of open half-planes, hence connected whenever nonempty.
Apply Lemma~\ref{descent} to $T=S$ intersected with each such sign region;
its span dimension is at most 3, so it contains at most four points.
We have therefore placed full $\nu$-measure in a finite set of size bounded by \eqref{bound}.
Its closure is the same finite set, so full support of $\nu$ implies that all of $S$
is contained in it, including any previously excluded boundary candidates.
Every support point is now an atom of positive mass. The individual zero-equator
statements in \eqref{all-rows} imply $v\cdot w\ne0$ for every pair of support atoms.
Taking a minimum over this finite set proves the asserted margin.
The actual fibers $F^{-1}(v)$ provide the partition, and \eqref{all-rows} identifies its rows.
\end{proof}
The margin asserted here may depend on the maximizer.

\section{Separating binary observations}\label{sec-separation}
By the preceding section, write $S=\{v_1,\ldots,v_m\}$ with actual masses
$p_i>0$. Put $V=(v_i^T)_{i=1}^m$, $P_0=\diag(p_i)$ and
$B_{ij}=\mathbf1_{\{v_i\cdot v_j>0\}}$. Then
\begin{equation}\label{finite-rows}
 BP_0V=VL,\qquad \sum_i p_iv_iv_i^T=I_3,\qquad
 v_i^TLv_i=\mathcal M.
\end{equation}
The full rows of $B$ are distinct: equal rows would imply equal $v_i$
by invertibility of $L$. A positive clique $J$ means a subset of this actual support with $B_{ij}=1$
for every $i,j\in J$. Extend vectors indexed by $J$ by zero outside $J$.
\begin{lemma}[Vector rigidity on a positive clique]
For every real $a$ supported on a positive clique, if
$Q_a=\sum_i a_i v_iv_i^T=0$, then
\begin{equation}
 \sum_i a_i=0,\quad b_a:=\sum_i a_iv_i=0,\quad
 \sum_i B_{ki}a_iv_i=0\quad\hbox{for every row }k. \label{rigid}
\end{equation}
\end{lemma}
\begin{proof}
Taking the trace against $L$ gives $\mathcal M\sum a_i=0$.
For both signs of sufficiently small nonzero $t$, all $p_i+ta_i$ are positive.
These are legitimate new fiber measures for the same binary template, realizable
on an atomless space. They are not independent choices of moment data.
The original coordinate functions have Gram matrix $I+tQ_a=I$ and compression
\[
 N(t)=L+t(Q_aL+LQ_a)+t^2b_ab_a^T=L+t^2b_ab_a^T.
\]
The last term uses the clique assumption, and both cross terms use \eqref{finite-rows}.
Their trace cannot exceed the global maximum $\mathcal M$; hence $b_a=0$ and $N(t)=L$.
Equality of this trial trace with the maximum forces its three-dimensional
space to be invariant. Explicitly, rotate one orthonormal coordinate into any
unit perpendicular function; the first derivative of the trace gives zero
off-space matrix element. This argument does not differentiate ordered eigenvalues.
The compression $L$ then gives the eigenfunction equations at the perturbed measure.
All fibers still have positive mass, so for every $k$,
$\sum_i(p_i+ta_i)B_{ki}v_i=Lv_k$.
Subtract \eqref{finite-rows} and divide by $t$ to obtain \eqref{rigid}.
\end{proof}
Moment-preserving cloning and the Ky Fan equality argument follow the
mechanism in~\cite[Claims 2.2--2.5]{KMMP}. The permissible size of $t$
may depend on the maximizer and the chosen direction.

\subsection{A separating observation lemma}
Minimal dependencies of rank-one powers have a broader theory; see
Li and Liu~\cite{LiLiu}. The following lemma uses the additional binary
separation hypothesis. Its elementary proof is included.
\begin{lemma}[Moment independence from separating binary observations]\label{binary}
Let $T=\{v_i:i\in J\}$ be a finite set of distinct points in $\R^3$ with
$v_{i1}>0$ and $v_i^TLv_i=\mathcal M>0$, where $L$ is positive definite.
Let $\mathcal B$ be a family of functions $b:J\to\{0,1\}$ containing the
constant-one function and separating every pair of distinct indices.
Suppose, for every real $a$ with $\sum_{i\in J}a_iv_iv_i^T=0$, that
\begin{equation}
 \sum_{i\in J}a_i b(i)v_i=0\quad\hbox{for every }b\in\mathcal B. \label{obs}
\end{equation}
Then the matrices $\{v_iv_i^T:i\in J\}$ are linearly independent.
\end{lemma}
\begin{proof}
Suppose instead that $I\subseteq J$ is a support-minimal dependent family
of these matrices, with relation $a_i\ne0$ for all $i\in I$.
Write $s=|I|$. Since $\dim\operatorname{Sym}_3=6$, minimality gives $s\le7$.
The ellipse equation gives $\sum_{I}a_i=0$, while the constant observation
in \eqref{obs} gives $\sum_I a_iv_i=0$.
With $w_i=(1,v_i)^T\in\R^4$, all three identities combine to give
\[
 \sum_{a_i>0}a_iw_iw_i^T=\sum_{a_i<0}(-a_i)w_iw_i^T.
\]
Both sides are positive semidefinite. The range of a positive weighted sum
of $w_iw_i^T$ is precisely the span of those $w_i$:
its kernel consists of vectors perpendicular to every such $w_i$.
Consequently the two spans coincide, and the affine hull $H$ of all $v_i$
in the circuit satisfies
\[
 1+\dim H=\dim\operatorname{span}\{w_i:i\in I\}
 \le\min(|\{a_i>0\}|,|\{a_i<0\}|)\le3.
\]
This is an elementary positive-Gram range argument, not a new general circuit bound.
Distinct points here cannot be proportional: equal ellipse radius forces
the proportionality factor to be $\pm1$, and the positive first coordinate
excludes minus one. Also an affine line meets this positive-definite ellipsoid
in at most two points, by its strictly quadratic restriction to the line.
If $\dim H\le1$, there are at most two points, whose rank-one matrices
are independent by nonproportionality, a contradiction.
It remains that $\dim H=2$.

If $0\in H$, this is a two-dimensional linear plane. Its rank-one matrices
span a space of dimension at most three, so minimality now gives $s\le4$.
For any observation $b$ nonconstant on $I$, \eqref{obs} and the total vector
relation give nonzero-coefficient linear relations on both nonempty sets
$I_1=\{i:b(i)=1\}$ and $I_0=\{i:b(i)=0\}$.
Nonproportionality forces each set to contain at least three points.
This would give $s\ge6$, a contradiction. If instead every observation in
$\mathcal B$ were constant on $I$, then minimality excludes $s=1$ (a
one-element circuit is impossible because $v_iv_i^T\ne0$), so we may choose
distinct $i,j\in I$; since $\mathcal B$ separates every pair of distinct
indices of $J$, some $b\in\mathcal B$ satisfies $b(i)\ne b(j)$,
contradicting constancy on $I$.

If $0\notin H$, write $H=\{v:c\cdot v=1\}$.
Any three distinct points of $H$ on the ellipsoid are linearly independent:
a linear relation has coefficient sum zero after applying $c$, hence would
make the three points affinely collinear, which the line-intersection fact forbids.
A nonconstant observation would therefore require at least four points
on each of $I_1,I_0$, yielding $s\ge8$, again a contradiction.
Thus every observation is constant on $I$. A one-element circuit is impossible
because $v_iv_i^T\ne0$; separation of any two indices in $I$ is contradicted.
\end{proof}

\begin{theorem}[Every actual positive clique has at most six types]\label{six}
For every actual maximizer, and every positive
clique $J$ of its full support, the map
$a\mapsto\sum_{i\in J}a_iv_iv_i^T$ is injective. In particular $|J|\le6$.
If $\{v_i:i\in J\}$ lies in a linear plane, then $|J|\le3$.
If it lies in an affine plane not containing zero, then $|J|\le5$.
\end{theorem}
\begin{proof}
Use the actual observations $b_k(i)=B_{ki}$, for all rows $k$.
For nonempty $J$, a row $k\in J$ supplies the constant-one observation.
They separate indices: $B$ is symmetric and has distinct full rows.
Lemma \ref{binary} now applies by \eqref{rigid}.
The linear-plane refinement follows from $\dim\operatorname{Sym}_2=3$.
For an affine plane $c\cdot v=1$, the nonzero symmetric matrix
$L-\mathcal M cc^T$ is orthogonal to every $v_iv_i^T$. It is nonzero because $L$
has rank three while $\mathcal M cc^T$ has rank at most one. This leaves dimension at most five.
Empty cliques require no argument.
\end{proof}

\begin{corollary}[Proof of Theorem~\ref{structure-main}]\label{twentyfour}
Every maximizer has at most 24 distinct canonical support points.
\end{corollary}
\begin{proof}
Partition its full support by the signs of coordinates two and three,
assigning a zero coordinate to the positive sign. There are at most four parts.
Within each part, the last two coordinate products are nonnegative and the
first-coordinate product is strictly positive. Each part is a positive clique.
Theorem \ref{six} gives $m\le4\cdot6=24$ on its original support. This proves Theorem~\ref{structure-main}.
\end{proof}

\section{Moments and regression for one isotropic law}\label{sec-moments}
We now turn to the numerical estimate. Throughout this section the law
is arbitrary and isotropic. Bounded positive radii will be imposed only
when a radial comparison requires them. For later bookkeeping define
\begin{equation}\label{constants}
 d_0=-\frac{25624284539519}{53655078400000}
                  +\frac{\sqrt5}{2}-\frac{41\sqrt3}{120},\qquad
 d_1=d_0+\frac{103}{9000000},\qquad
 d=d_1+\frac1{3000000}.
\end{equation}
Then $d=(3+\sqrt5)/2-2c$, with $c$ as in~\eqref{main-constant}.
The intermediate constants separate the joint moment estimate from
the final atomic term.
Fix $\E XX^T=I_3$, let $R=|X|$, and put $U=X/R$ when $R>0$.
At $R=0$ choose any fixed unit vector for $U$. Define
\[
 \begin{gathered}
 C=\E R,\quad V=3-C^2,\quad \mathbf m=\E X,\quad B=|\mathbf m|^2,\quad e=1-B,\\
 A=\E|X\cdot X'|,\qquad D=\frac{3+\sqrt5}{2}-A-B.
 \end{gathered}
\]
Here $V,e\ge0$: $\E R^2=3$ and Cauchy--Schwarz gives
$(v\cdot \mathbf m)^2\le\E(v\cdot X)^2=1$ for every unit $v$.
All pair expectations below are over independent copies of this same law.
Set
\[
 M=\E[R UU^T]=\frac C3I_3+Z,\qquad
 J=1-\tr(M^2)=\frac V3-\|Z\|_F^2.
\]
The regression vector $\xi=U-MX$ satisfies
\begin{equation}\label{regression}
 \E[X\xi^T]=0,\quad \E\|\xi\|^2=J\ge0,\quad
 \E UU^T=M^2+\E[\xi\xi^T].
\end{equation}
In particular $J$ is an actual regression error, not a freely chosen moment.

Here are the harmonic normalizations used throughout. All tensor inner
products sum over all ordered indices. Put
\[
 \Omega_{ijkl}=\frac{\delta_{ij}\delta_{kl}+\delta_{ik}\delta_{jl}
                           +\delta_{il}\delta_{jk}}{15},
 \qquad
 ({\cal L}_3(y))_{ijk}
       =\frac{y_i\delta_{jk}+y_j\delta_{ik}+y_k\delta_{ij}}5.
\]
For symmetric traceless $Z$, put
\[
 ({\cal L}_4(Z))_{ijkl}=\frac17(\delta_{ij}Z_{kl}+\delta_{ik}Z_{jl}
 +\delta_{il}Z_{jk}+\delta_{jk}Z_{il}+\delta_{jl}Z_{ik}+\delta_{kl}Z_{ij}).
\]
For unit $u$, set
\[
 \begin{gathered}
 \Psi_3(u)=u^{\otimes3}-{\cal L}_3(u),\qquad
 \Psi_4(u)=u^{\otimes4}-\Omega-{\cal L}_4(uu^T-I_3/3),\\
 T_i=\E[R\Psi_i(U)],\qquad H=\|T_3\|_F,\qquad K=\|T_4\|_F.
 \end{gathered}
\]
Contraction in two indices gives the orthogonal decompositions
\[
 \E[R U^{\otimes3}]=T_3+{\cal L}_3(\mathbf m),\qquad
 \E[R U^{\otimes4}]=T_4+C\Omega+{\cal L}_4(Z).
\]
Indeed $\operatorname{ct}\Omega=I_3/3$, $\|\Omega\|^2=1/5$,
$\operatorname{ct}{\cal L}_4(Z)=Z$, and
$\langle T,{\cal L}_4(Z)\rangle=(6/7)\langle\operatorname{ct}T,Z\rangle$.
Also $\operatorname{ct}{\cal L}_3(y)=y$ and
$\langle N,{\cal L}_3(y)\rangle=(3/5)\operatorname{ct}N\cdot y$.
It follows directly that
\begin{align}\label{harmonics}
 \langle\Psi_3(u),\Psi_3(v)\rangle&=g_3(u\cdot v),&
 g_3(t)&=t^3-3t/5,&\|\Psi_3(v)\|^2&=2/5,\notag\\
 \langle\Psi_4(u),\Psi_4(v)\rangle&=g_4(u\cdot v),&
 g_4(t)&=t^4-6t^2/7+3/35,&\|\Psi_4(v)\|^2&=8/35.
\end{align}
The tensor moments involve $R$, not $R^3$ or $R^4$, and are integrable
under $\E R^2=3$.

Write $s=\sqrt5$ and define the scalar absolute-value majorant and its slack by
\[
 a=\frac{1+s}{6},\quad b=\frac{7s-5}{8},\quad\rho=\frac{15s-25}{16},
 \qquad h_0(z)=\frac{(1-z)(sz-1)^2(sz+s+2)}{2(s+1)^2}.
\]
Expansion gives, for $|t|\le1$,
\[
 a+b(t^2-1/3)-\rho(t^4-1/5)-|t|=h_0(|t|)\ge0.
\]
This majorant is from~\cite[Lemma 4.4]{STW}.
Let $S=\E[RR'h_0(|U\cdot U'|)]$. The decompositions above give
\[
 A=aC^2+b\|Z\|^2-\rho\left(K^2+\frac67\|Z\|^2\right)-S.
\]
Consequently
\begin{equation}\label{budget4}
 D=e+\alpha V+\beta J+\rho K^2+S,
 \qquad \alpha=\frac{2s-1}{14},\quad\beta=\frac{s+10}{14}.
\end{equation}
Useful rational bounds are $\alpha\ge26/105$, $\beta\ge6/7$,
and $\rho\ge1/2$, using $s>67/30$, $s\ge2$, and $s\ge11/5$.

\subsection{A regression estimate}
Suppose $B>0$. Put
\[
 x=\sqrt B,\quad v=\mathbf m/x,\quad t=U\cdot v,\quad Y=Rt,\quad
 q_0=1/\sqrt3,\quad q=C/(3x),\quad L=\E[R(t-q_0)^2].
\]
Then $\E Y=x$, $\E Y^2=1$, $\operatorname{Var}Y=e$.
If $w_0=\E\xi$, \eqref{regression} gives
$w_0=\E[(1-\mathbf m\cdot X)\xi]$, hence $|w_0|^2\le eJ$.
For traceless symmetric $3$ by $3$ matrices,
$|Zv|^2\le(2/3)\|Z\|^2$; the other two eigenvalues accompanying an
eigenvalue $z$ have squared sum at least $z^2/2$.
Writing $J_v=\E(v\cdot\xi)^2\le J$, we obtain
\begin{align*}
 \E(t-q)^2
 &=|Mv-q\mathbf m|^2+q^2e+J_v-2qv\cdot w_0\\
 &\le 2V/9+J/3+q^2e+2q\sqrt{eJ}=:F_0,\\
 \E(Y-qR)^2&\le V/3+3q^2e+2q\sqrt{eV}=:G_0.
\end{align*}
Here $Mv-q\mathbf m=Zv$ and $\E RY\ge Cx-\sqrt{eV}$.
Cauchy--Schwarz gives $L_q:=\E[R(t-q)^2]\le\sqrt{F_0G_0}$.
Note that $Rt=R(U\cdot v)=(RU)\cdot v=X\cdot v=Y$, so $\E[Rt]=\E Y=x$;
this is the only place at which the regression vector enters the exact
identity below. With $\zeta=C/\sqrt3$,
\[
 L_q-L=q_0(\zeta-x)^2(\zeta+2x)/x^2\ge0.
\]
Thus, writing $u=\sqrt e$, $z=\sqrt V$, $j=\sqrt J$,
\begin{equation}\label{Lbar}
 L\le\overline L:=\frac34F_0+\frac13G_0
 =\frac{7q^2}{4}u^2+\frac5{18}z^2+\frac14j^2
                       +\frac{2q}{3}uz+\frac{3q}{2}uj.
\end{equation}
Set $r_{\rm reg}=\overline L-L\ge0$. This remainder will be kept
in the joint identity. No lower radius bound has been used here.

\section{A signed quartic and a joint defect identity}\label{sec-quartic}
Define
\[
 f(w)=\frac{8352+43806w+54016w^2-12928w^3-10240w^4}{41503},
 \qquad \sigma(w)=f(w)-|w|-w.
\]
The polynomial is the degree-at-most-four Hermite interpolant with double
contacts to $|w|+w$ at $-3/8$ and $1/2$, and a contact at $1$.
Uniqueness follows because a difference of two such polynomials would have
five zeros counted with multiplicity. The following explicit identities prove
the needed inequalities; interpolation alone is not a majorant proof.
\begin{lemma}\label{scalar}
For every $w\in[-1,1]$,
\begin{equation}\label{sigmapoint}
 \sigma(w)\ge0,\qquad \sigma(w)\ge\frac{1-9w^2}{25}.
\end{equation}
\end{lemma}
\begin{proof}
Direct expansion gives
\[
 41503\sigma(w)=
 \begin{cases}
 (1-w)(2w-1)^2(8352+2560w),&0\le w\le1,\\
 (8w+3)^2(928-82w-160w^2),&-1\le w\le0.
 \end{cases}
\]
The final factor on the negative interval is concave and at least $850$.
This proves nonnegativity, including the join at $0$ and both endpoints.
The second assertion follows immediately when $|w|\ge1/3$.
For $w=-z$, $0\le z\le1/3$, write
\[
 \sigma(-z)=(z-3/8)^2
          \frac{59392+5248z-10240z^2}{41503}
       \ge\frac{10}{7}(z-3/8)^2.
\]
Indeed $z(5248-10240z)\ge0$ on this interval and
$59392/41503>10/7$. The remaining comparison is the exact square
\[
 \frac{10}{7}(z-3/8)^2-\frac{1-9z^2}{25}
  =\frac{313}{175}\left(z-\frac{375}{1252}\right)^2
                         +\frac{109}{250400}>0.
\]
For $w=z$, $0\le z\le1/3$, the positive-interval factorization gives
\[
 \sigma(z)=(z-1/2)^2\frac{(1-z)(33408+10240z)}{41503}
              \ge\frac12(z-1/2)^2,
\]
because the coefficient is at least $22272/41503>1/2$. Now use
\[
 \frac12(z-1/2)^2-\frac{1-9z^2}{25}
   =\frac{43}{50}\left(z-\frac{25}{86}\right)^2+\frac{53}{4300}>0.
\]
These intervals exhaust $[-1,1]$.
\end{proof}

Put $P=\E[RR'\sigma(U\cdot U')]$ for the same actual law as before, and set
\[
 \begin{gathered}
 a_3=\frac{12928}{41503},\quad a_4=\frac{10240}{41503},\quad
 b_e=\frac{16386}{18865},\\
 b_V=\frac{64608}{290521},\quad b_J=\frac{316672}{290521},\quad
 \delta_{\mathrm q}=\frac{3+\sqrt5}{2}-\frac{544886}{207515}.
 \end{gathered}
\]
Integration of the expanded $f$, using \eqref{harmonics}, gives the single
exact joint budget
\begin{equation}\label{budget7}
 D=\delta_{\mathrm q}+b_e e+b_V V+b_J J+a_3H^2+a_4K^2+P.
\end{equation}
For example, writing $f=c_0+c_1w+c_2w^2-a_3w^3-a_4w^4$,
the coefficients are $b_e=c_1-3a_3/5$, $b_V=c_0+3a_4/35$,
$b_J=c_2-6a_4/7$, and the integrated constant is
$3(c_0-a_4/5)+b_e+c_2=544886/207515$.
All tensors are $R$-weighted and $\sigma$ is bounded, so this equality
uses no radial moment beyond the second one. Identity~\eqref{budget4} will only justify a small-$D$ branch; it is
not added as an independent resource to~\eqref{budget7}.

\begin{lemma}\label{radialpayment}
Suppose $0<r_-\le R\le r_+<\infty$ almost surely and $r_+/r_-\le3$.
Then
\begin{equation}\label{payment}
 P=\frac3{50}(V-9J)+r_{\rm rad},\qquad r_{\rm rad}\ge0.
\end{equation}
\end{lemma}
\begin{proof}
For $w=U\cdot U'$, the exact pair identities from the full Gram are
\[
 V=\frac12\E(R-R')^2,\qquad
 J=\frac12\E[(R-R')^2w^2].
\]
The latter uses $\E[R^2w^2]=\E[(R')^2w^2]=1$ and
$\E[RR'w^2]=\tr(M^2)$; it assumes no independence of radius and direction.
The ratio bound gives $(R-R')^2\le(4/3)RR'$. Use the two conclusions
of Lemma~\ref{scalar} in this order:
\[
 P\ge\frac34\E[(R-R')^2\sigma(w)]
   \ge\frac3{100}\E[(R-R')^2(1-9w^2)]
    =\frac3{50}(V-9J).
\]
The expression $V-9J$ need not be nonnegative. Define $r_{\rm rad}$
as the actual difference; only this difference is claimed nonnegative.
\end{proof}

\section{Directional and radial control}\label{sec-directional}
Throughout this section assume $0<r_-\le R\le r_+<\infty$ almost surely,
with $r_+/r_-\le3$, and work in the branch $D<1/20$.
Equation~\eqref{budget4} implies
\begin{equation}\label{smallbranch}
 e<1/20,\quad x>9/10,\quad V<21/104<2/9,\quad C>5/3,
 \qquad q^2=\frac{3-V}{9(1-e)}\le\frac{20}{57}.
\end{equation}
In particular $\mathbf m\ne0$. Put $p_i=\langle T_i,\Psi_i(v)\rangle
=\E[Rg_i(t)]$. The elementary Hermite identities
\begin{align}\label{Hermite}
 \left(\frac{53}{28}t^2-\frac94t+\frac{41}{70}\right)-g_4(t)
       &=(1-t)(t-1/2)^2(t+2)\ge0,\notag\\
 (2t^2-37t/20+1/4)-g_3(t)
       &=(1-t)(t-1/2)^2\ge0
\end{align}
hold on $[-1,1]$. Define
\[
 \gamma_4=9/4-53\sqrt3/42,\quad
 \ell_4=9/4-511\sqrt3/420,\quad
 \ell_3=37/20-11\sqrt3/12.
\]
Here $0<\gamma_4<1/15$ (square the positive rational comparisons).
Since $\E Rt=x$, $\E Rt^2=L+2x/\sqrt3-C/3$, the first identity gives
\[
 -p_4\ge\gamma_4x+19C/420-53L/28
       \ge\ell_4-e/15-19V/1400-53L/28.
\]
We used $x\ge1-e$ and $C\ge\sqrt3-3V/10$; the latter follows from
$V=(\sqrt3-C)(\sqrt3+C)$ and $C>5/3$.
The second identity gives
\[
 -p_3\ge(37/20-4/\sqrt3)x+5C/12-2L
       \ge\ell_3+9e/40-V/8-2L.
\]
Here $4/\sqrt3-37/20>9/20$ and $x\le1-e/2$.
Define the nonnegative differences in these last two inequalities by
\begin{align*}
 r_4&=-p_4-\ell_4+e/15+19V/1400+53L/28,\\
 r_3&=-p_3-\ell_3-9e/40+V/8+2L.
\end{align*}
The polynomial identities, rather than a fixed sphere assumption,
justify these directional estimates for the actual weighted moments.

Set $\eta=1/10$, $\theta=6/25$, and
$\ell=53\eta/28+2\theta=937/1400$. Define the tensor squares
\[
 {\cal R}_4=a_4\left\|T_4+\frac{\eta}{2a_4}\Psi_4(v)\right\|^2,\qquad
 {\cal R}_3=a_3\left\|T_3+\frac{\theta}{2a_3}\Psi_3(v)\right\|^2.
\]
Their exact expansions give
\[
 a_4K^2=-\eta p_4-\frac{2\eta^2}{35a_4}+{\cal R}_4,\qquad
 a_3H^2=-\theta p_3-\frac{\theta^2}{10a_3}+{\cal R}_3.
\]
Define
\[
 E=\frac{5183683}{5659500},\quad
 V_c=\frac{145858703}{581042000},\quad
 J_c=\frac{7989533}{14526050},\quad
 {\cal R}_0={\cal R}_3+{\cal R}_4+\eta r_4+\theta r_3\ge0.
\]
Substituting the directional identities and \eqref{payment} into
\eqref{budget7}, then writing $L=\overline L-r_{\rm reg}$, gives the exact identity
\begin{equation}\label{certificate-start}
 D=d_0+E e+V_c V+J_c J-\ell\overline L
             +{\cal R}_0+r_{\rm rad}+\ell r_{\rm reg}.
\end{equation}
Here
\[
 d_0=\delta_{\mathrm q}+\eta\ell_4+\theta\ell_3
                 -\frac{2\eta^2}{35a_4}-\frac{\theta^2}{10a_3},
\]
which simplifies to \eqref{constants}. The coefficient checks are
\[
 E=b_e-\eta/15+9\theta/40,\quad
 V_c=b_V-19\eta/1400-\theta/8+3/50,\quad J_c=b_J-27/50.
\]
All the remainders in this identity refer to the same probability law.

Put $B_c=V_c-5\ell/18$ and $C_c=J_c-\ell/4$. The exact rational values are
\[
 B_c=\frac{42565069}{653672250},\qquad
 C_c=\frac{88944217}{232416800}.
\]
For $u=\sqrt e$, $z=\sqrt V$, and $j=\sqrt J$, the remaining quadratic has
the explicit completion
\begin{align}\label{squares}
 E e+V_c V+J_c J-\ell\overline L
 &=\frac{13}{200}\left(z-\frac{937}{273}qu\right)^2
   +\frac{19}{50}\left(j-\frac{2811}{2128}qu\right)^2\notag\\
 &\quad+\left(\frac{183}{200}-\frac{725029049}{278853120}q^2\right)e\notag\\
 &\quad+\frac{10481}{11319000}e
        +\frac{305491}{2614689000}V+\frac{625833}{232416800}J.
\end{align}
Indeed the last three fractions are $E-183/200$, $B_c-13/200$,
and $C_c-19/50$. Also
\[
 \frac{725029049}{278853120}
  =\frac{7\ell}{4}+\frac{\ell^2}{9(13/200)}
                          +\frac{9\ell^2}{16(19/50)}.
\]
By \eqref{smallbranch}, the coefficient of $e$ in parentheses is at least
\[
 \frac{183}{200}-\frac{20}{57}\frac{725029049}{278853120}
                         =\frac{53754367}{19868284800}>0.
\]
Thus every term on the right of \eqref{certificate-start} after $d_0$
is nonnegative. The displayed squares give an explicit certificate for the remaining
quadratic form.

\begin{proposition}\label{conditional-law}
Every isotropic law on $\R^3$ whose radii lie almost surely in a positive
interval of endpoint ratio at most $3$ satisfies $D\ge d_0$.
In the branch $D<1/20$, \eqref{certificate-start}--\eqref{squares}
give the full common-remainder certificate.
\end{proposition}
\begin{proof}
The small branch was proved above. Otherwise $D\ge1/20>d_0$.
For explicit exact comparisons use
$682/305<\sqrt5<161/72$ and $433/250<\sqrt3<97/56$,
each verified by squaring. They yield
\[
 d_0-7/144>\frac{940735364069}{29456638041600000}>0,\qquad
 1/20-d_0>\frac{620701786231}{482895705600000}>0.
\]
No $B=0$ division occurs: it is excluded by $e<1/20$ in the small branch.
\end{proof}

\section{A reserve in the quartic slack}\label{sec-reserve}
Throughout this section assume $0<r_-\le R\le r_+<\infty$ almost surely,
with $r_+/r_-\le\sqrt8$. Write $g(w)=(1-9w^2)/25$. Then
\[
 \frac{(R-R')^2}{RR'}\le\sqrt8+\frac1{\sqrt8}-2
    =\frac{9\sqrt2}{4}-2\le\frac{32}{27}.
\]
For the last inequality, $\sqrt2<344/243$ follows by squaring, with
$(344/243)^2-2=238/59049>0$. Put
$z_{\rm rad}=3(R-R')^2/(4RR')$, so $0\le z_{\rm rad}\le8/9$.
The actual remainder in \eqref{payment} now has the exact decomposition
\begin{equation}\label{reserve}
 r_{\rm rad}=\frac P9+r_{\rm tr},\qquad
 r_{\rm tr}=\E\!\left[RR'\left\{(8/9-z_{\rm rad})\sigma(w)
                    +z_{\rm rad}(\sigma(w)-g(w))\right\}\right]\ge0.
\end{equation}
This uses both pointwise inequalities of Lemma~\ref{scalar}, including
where $g<0$. It is not a second independent budget or a new copy of $P$.

\subsection{A contact polynomial without a cubic term}
Using the same interpolation nodes and tensor weights, set
\[
 h_*(t)=(t-1)(t-1/2)(t+3/8),\qquad
 z_*(t)=(t+9/8)h_*(t)
       =t^4-\frac{85}{64}t^2+\frac{15}{128}t+\frac{27}{128}.
\]
The factor $t+9/8$ cancels the cubic term. We need the following elementary
bound, for which we give an explicit proof.
\begin{lemma}\label{contact-square}
For $-1\le t\le1$, $z_*(t)^2\le3\sigma(t)/2$.
\end{lemma}
\begin{proof}
For $0\le t\le1$, cancel common factors away from their zeros to obtain
\[
 \frac{z_*(t)^2}{\sigma(t)}
  =\frac{41503}{4(8352+2560t)}q_*(t),\qquad
 q_*(t)=(1-t)(t+3/8)^2(t+9/8)^2.
\]
The cancelled expressions extend continuously at the zeros. Its derivative is
\[
 q_*'(t)=(t+3/8)(t+9/8)(165/64-t/2-5t^2).
\]
The last factor decreases on $[0,1]$, equals $13/576>0$ at $2/3$,
and equals $-33/256<0$ at $11/16$. Thus the maximum is within
$[2/3,11/16]$. On that interval $(t+3/8)(t+9/8)<2$ and
$q_*'(t)\le13/288$. The interval has length $1/48$, so
\[
 q_*(t)\le\frac{1155625}{995328}+\frac{13}{13824}<\frac65;
\]
the final positive difference is $189163/4976640$.
Consequently $z_*^2/\sigma\le41503/27840<3/2$ on this half interval.
For $t=-x$, $0\le x\le1$, the other factorization gives
\[
 \frac{z_*(-x)^2}{\sigma(-x)}
 =\frac{41503\,p_*(x)^2}{64(928+82x-160x^2)},\qquad
 p_*(x)=(1+x)(x+1/2)(9/8-x)\ge0.
\]
The quadratic denominator factor is at least $850$, by concavity and its
endpoint values. With $y=x-1/2\in[-1/2,1/2]$,
\[
 p_*(x)=15/16+y/16-15y^2/8-y^3
 \le15/16+y/16-11y^2/8
 =\frac{1321}{1408}-\frac{11}{8}(y-1/44)^2<1.
\]
It follows that $z_*^2/\sigma\le41503/54400<3/2$ here.
At every cancelled zero the claimed inequality follows directly or by
continuity. These two intervals cover the domain.
\end{proof}

\begin{remark}\label{rem-theta}
The constant $3/2$ is not sharp: a dense grid scan gives
$\sup_{[0,1]}z_*^2/\sigma=1.200224\ldots$ and
$\sup_{[-1,0]}z_*^2/\sigma=0.615384\ldots$. We keep $3/2$ because it keeps
every subsequent coefficient rational and hand-verifiable. The loss enters
the certificate only through $r_{\rm ct}$ and through the divisor $125$ in
\eqref{Cstar8}; both are quantified in Appendix~\ref{sec-appendix}.
\end{remark}

In particular the following displayed square proves a useful linear payment:
\begin{equation}\label{contact-paid}
 \sigma(t)-\frac{z_*(t)}{125}+\frac3{125000}
   =\left(\sigma(t)-\frac23z_*(t)^2\right)
       +\frac23\left(z_*(t)-\frac3{500}\right)^2\ge0.
\end{equation}
All these functions are bounded on $[-1,1]$. Integrating against the actual
nonnegative weight $RR'$ uses only $\E R<\infty$, not a fourth or eighth
radial moment. The actual contact moment is
\begin{align}\label{Zstar}
 Z_*:=\E[RR'z_*(U\cdot U')]
   &=K^2+\frac7{320}-\frac{15}{128}e
                    -\frac{561}{4480}V+\frac{211}{448}J.
\end{align}
Indeed the $RR'$ moments of $1,w,w^2,w^4$ are respectively
$3-V,1-e,1-J,K^2+6\|Z\|^2/7+(3-V)/5$, all for this same law.
Define
\begin{equation}\label{Cstar8}
 {\cal C}=Z_*/125-3C^2/125000,\qquad
 r_{\rm ct}=P-{\cal C}\ge0.
\end{equation}
The scalar ${\cal C}$ need not itself be nonnegative. Its exact contribution is
\begin{equation}\label{paid-expansion}
 \frac{{\cal C}}{9}=\frac{103}{9000000}+\frac{K^2}{1125}
 -\frac e{9600}+\left(-\frac{187}{1680000}+\frac1{375000}\right)V
                       +\frac{211}{504000}J.
\end{equation}

\subsection{The full joint certificate}
For $D<1/20$ define the explicitly nonnegative quadratic expression
\begin{align}\label{Q8}
 Q={}&\frac{13}{200}\left(\sqrt V-\frac{937}{273}q\sqrt e\right)^2
       +\frac{19}{50}\left(\sqrt J-\frac{2811}{2128}q\sqrt e\right)^2\notag\\
 &+\left(\frac{183}{200}-\frac{725029049}{278853120}q^2\right)e\notag\\
 &+\frac{148831}{181104000}e
       +\frac{42847433}{5229378000000}V+\frac{65082103}{20917512000}J.
\end{align}
The parenthetical coefficient is positive by \eqref{smallbranch} and the
same exact lower bound following \eqref{squares}. The last coefficients are
obtained from those in \eqref{squares}, respectively by adding
$-1/9600$, $-187/1680000+1/375000$, and $211/504000$.
Substitution of \eqref{reserve}--\eqref{paid-expansion} into \eqref{certificate-start} yields the exact identity
\begin{equation}\label{certificate8}
 D=d_1+Q+{\cal R}_0+\ell r_{\rm reg}
          +\frac{K^2}{1125}+r_{\rm tr}+\frac{r_{\rm ct}}9.
\end{equation}
Every remainder is defined from the same $X$, its independent copy and its
complete Gram matrix.
\begin{proposition}\label{conditional8}
Every isotropic law whose radii lie almost surely in a positive bounded
interval of endpoint ratio at most $\sqrt8$ satisfies $D\ge d_1$.
In the branch $D<1/20$, \eqref{certificate8} retains all displayed remainders.
\end{proposition}
\begin{proof}
In the small branch all terms after $d_1$ are nonnegative. In the other
branch $D\ge1/20>d_1$. To verify the last comparison entirely by rationals,
use $\sqrt5<161/72$ and $\sqrt3>433/250$ in \eqref{constants}; they give
\[
 \frac1{20}-d_1>\frac{1025292188593}{804826176000000}>0.
\]
No division by $B=0$ occurs in the small branch, and the other branch
requires none. The radial hypothesis will be verified for a maximizing law in
Section~\ref{sec-conclusion}, rather than imposed on every isotropic law.
\end{proof}

\section{The coefficient bound and its transfer}\label{sec-conclusion}
We return to a maximizing graphon. Put
\[
 D_*=(3+\sqrt5)/2-2\mathcal M.
\]
By the threshold identity in Lemma~\ref{contact-structure}, its canonical
law $X=F$ satisfies $A+B=2\mathcal M$, so the deficit $D$ of the preceding
sections equals $D_*$.
\begin{lemma}\label{spectral-reserve}
If $W$ is an actual maximizer and $D_*=(3+\sqrt5)/2-2{\cal M}<1/20$,
then $\lambda_3>44/375$, $\lambda_1<15/16$, and the canonical frame
has positive bounded radii of endpoint ratio less than $\sqrt8$.
\end{lemma}
\begin{proof}
The small branch and $\sqrt5>682/305>559/250$ give
\[
 {\cal M}>\frac{3+\sqrt5}{4}-\frac1{40}>\frac{321}{250}.
\]
Lemma~\ref{ranktwo} gives
$\lambda_3={\cal M}-\Sigma_2(W)>321/250-7/6=44/375$.
For $a=\lambda_1$, the Hilbert--Schmidt trace and the constant-function trial give
\[
 \lambda_1^2+\lambda_2^2+\lambda_3^2\le\iint W^2\le\iint W\le\lambda_1,
 \qquad {\cal M}\le a+\sqrt{2a(1-a)},\quad 0<a\le1.
\]
The function on the right is decreasing on $[15/16,1]$; differentiation
on the interior reduces this to $(2a-1)^2>2a(1-a)$ there.
Moreover
\[
 \frac{15}{16}+\sqrt{\frac{15}{128}}<\frac{321}{250},\qquad
 \left(\frac{321}{250}-\frac{15}{16}\right)^2-\frac{15}{128}
       =\frac{11499}{4000000}>0.
\]
Thus $a<15/16$. By the actual ellipsoid identity,
\[
 \sqrt{{\cal M}/\lambda_1}\le R\le\sqrt{{\cal M}/\lambda_3},\qquad
 \frac{\esssup R}{\essinf R}
 \le\sqrt{\lambda_1/\lambda_3}<\sqrt{5625/704}<\sqrt8,
\]
where $8-5625/704=7/704$. No simplicity assumption was made on
the positive eigenvalues; the argument uses the whole orthonormal frame.
\end{proof}

\begin{remark}\label{rem-radii}
Lemma~\ref{radialpayment} is stated with endpoint ratio $3$, whereas the
present lemma supplies $\sqrt8<3$ for an actual maximizer; the two
hypotheses are therefore compatible, with slack. Using $\sqrt8$ in
Lemma~\ref{radialpayment} would replace the factor $4/3$ in \eqref{payment}
by $32/27$, a $12.5\%$ gain on $P$, which is not needed here.
\end{remark}

\subsection{Positive-mass self-pairs}
By Theorem~\ref{structure-main}, write the canonical law as
$\sum_{i=1}^m p_i\delta_{v_i}$ with $m\le24$ and all $p_i>0$.
Set $R_i=|v_i|$, $u_i=v_i/R_i$, $w_{ij}=u_i\cdot u_j$, and
\[
 \kappa=\frac3{125000},\qquad
 h(t)=\sigma(t)-\frac{z_*(t)}{125}+\kappa,
 \qquad \chi=\sum_i p_i^2R_i^2.
\]
Equation~\eqref{contact-paid} gives $h\ge0$ and $h(1)=\kappa$.
The same contact remainder in~\eqref{Cstar8} therefore has the exact split
\begin{equation}\label{atom-split}
 r_{\rm ct}=\kappa\chi+r_{\rm off},\qquad
 r_{\rm off}=\sum_{i\ne j}p_ip_jR_iR_jh(w_{ij})\ge0.
\end{equation}
These self-pairs have positive product mass. They are not values assigned
to the null diagonal $x=y$ of the underlying nonatomic space.
The estimate $\chi\ge C^2/m$ is the usual Cauchy--Schwarz
cardinality bound. For related weighted-energy arguments,
see~\cite[Theorem 2.1 and the subsequent discussion]{Weighted}.

In the branch $D_*<1/20$, let
\[
 \widehat Q=Q-\frac{V}{9000000}.
\]
It is nonnegative: in~\eqref{Q8} only the last separate $V$ coefficient
changes, to
\[
 \frac{42847433}{5229378000000}-\frac1{9000000}
       =\frac{14088797}{1743126000000}>0.
\]
\begin{proposition}\label{final-certificate}
For every maximizing graphon with $D_*<1/20$,
\begin{align}\label{certificate-final}
 D_*={}&d+\widehat Q+{\cal R}_0+\ell r_{\rm reg}
       +\frac{K^2}{1125}+r_{\rm tr}\notag\\
 &+\frac{\kappa}{9}\left(\chi-\frac{C^2}{24}\right)
       +\frac{r_{\rm off}}9.
\end{align}
Every term after $d$ is nonnegative. Consequently $\mathcal M\le c$.
\end{proposition}
\begin{proof}
Lemma~\ref{spectral-reserve} verifies the radial hypothesis of
Proposition~\ref{conditional8}. Substitute~\eqref{atom-split} once into
\eqref{certificate8}. Since $C^2=3-V$,
\[
 \frac{\kappa\chi}{9}
 =\frac1{3000000}-\frac{V}{9000000}
       +\frac{\kappa}{9}\left(\chi-\frac{C^2}{24}\right).
\]
This gives~\eqref{certificate-final}. Moreover
\[
 \chi-\frac{C^2}{24}
 =\sum_i\left(p_iR_i-\frac Cm\right)^2
       +C^2\left(\frac1m-\frac1{24}\right)\ge0.
\]
No lower bound on an individual mass is needed. The complementary branch
also gives $D_*\ge d$, since the rational estimates used above imply
\[
 \frac1{20}-d>
       \frac{1025023913201}{804826176000000}>0.
\]
Thus $D_*\ge d$ in both branches, or equivalently $\mathcal M\le c$.
\end{proof}
The atomic term changes the coefficient by $1/6000000$ relative to
$(3+\sqrt5)/4-d_1/2$. Most of the improvement in Theorem~\ref{main}
comes from the preceding joint moment estimate, not from this last term.

\subsection{All isotropic laws and the exact finite shift}
\begin{proof}[Proof of Theorem~\ref{main}]
Any Borel law on $\R^3$ can be realized on a standard nonatomic probability
space, by adjoining an independent uniform variable if necessary.
For its isotropic $X$, the coordinates are an orthonormal three-frame and
$W_X=\mathbf1_{\{X\cdot X'>0\}}$ is a measurable symmetric binary kernel.
Therefore
\[
 \frac{\E|X\cdot X'|+|\E X|^2}{2}
   =\iint W_X X\cdot X'
   \le\Sigma_3(W_X)\le\mathcal M\le c.
\]
This proves~\eqref{law-bound}, including zero vectors and unbounded radii.
The radial restriction was used only for a maximizing canonical law.

For a finite graph, the exact-rank Ky Fan formula reads
\[
 \Lambda_3(G)=\max_{Q=Q^T=Q^2,\ \tr Q=3}\tr(A_GQ).
\]
Factor $Q=YY^T$ with orthonormal columns, let $y_i$ be its rows, and let
$X$ be uniform on $\sqrt n\,y_1,\ldots,\sqrt n\,y_n$. Then
\[
 \E XX^T=I_3,\qquad
 \E|X\cdot X'|=\|Q\|_1/n,\qquad
 |\E X|^2=\boldsymbol1^TQ\boldsymbol1/n,
\]
where $\|Q\|_1$ denotes the entrywise absolute sum.
Since $Q_{ii}\ge0$ and $\sum_iQ_{ii}=3$, the zero adjacency diagonal gives
\[
 \tr(A_GQ)\le\sum_{i\ne j}(Q_{ij})_+
 =\frac{\|Q\|_1+\boldsymbol1^TQ\boldsymbol1}{2}-3
 \le cn-3.
\]
Taking the maximum proves~\eqref{main-bound}. No positivity assumption on
$\lambda_3(G)$ is used; repeated eigenvalues, zero rows and $n=3$ cause
no exception.
\end{proof}

\section{Further questions}
The explicit construction in Lemma~\ref{attainment} leaves the interval
\[
 \frac{9+\sqrt5+2\sqrt{35}}{18}\le\mathcal M
       \le1.284673255723918505\ldots.
\]
The lower endpoint comes from a weighted wheel, but its global optimality
is not established here. The support bound reduces the structure of every
maximizer to finitely many positive-mass types without classifying the
possible templates. It does not imply that the binary matrix $B$ is
nonsingular, that every maximizer has six types, or that individual masses
are uniformly bounded away from zero. Determining the optimum remains
a separate problem.

The exact identity~\eqref{certificate-final} also distinguishes a
quantitative certificate from a qualitative equality obstruction. Any
further decrease obtained from this certificate must bound a remaining term from below
under the same constraints, without counting the same quartic remainder
twice. The present argument does not use an unquantified stability gap.

\appendix
\section{Constants: provenance and machine verification}\label{sec-appendix}

A reader cannot tell from the statements alone which of the many integer
coefficients in this paper are forced by an interpolation problem and which
were chosen and then verified. We record the distinction.

\textbf{Uniquely determined.}
(i)~The polynomial $f$ of \S6 is the unique polynomial of degree at most
four with double contact to $|w|+w$ at $w=-3/8$ and at $w=1/2$ and a contact
at $w=1$; five conditions determine five coefficients, so once the nodes are
fixed the integers in $f$ (including $41503=7^3\cdot11^2$) are forced.
(ii)~$z_*$ of \S8 is determined by the same nodes together with the
requirement that the cubic term vanish, i.e.\ by the factor $(t+9/8)$.
(iii)~$\cstar$ of Lemma~\ref{attainment} is the value of an explicit
one-variable function at the critical point $b=(7+5\sqrt7)/126$.
(iv)~The four rational certificates of \S7 and \S9 are obtained by
substituting the stated rational bounds for $\sqrt5$ and $\sqrt3$ into
$d_0,d_1,d$; no search is involved.

\textbf{Chosen, then verified exactly.}
(i)~The nodes $-3/8,\,1/2,\,1$ themselves, and the majorant $h_0$ taken
from~\cite[Lemma 4.4]{STW}.
(ii)~The parameters $\ell=937/1400$, $\eta=1/10$, $\theta=6/25$ of \S7 and
the square coefficients of \eqref{squares} and \eqref{Q8}; these were
selected so that all residual coefficients remain positive, and every
positivity claim is then checked by exact rational arithmetic.
(iii)~The divisor $125$ and $\kappa=3/125000$ of \eqref{Cstar8}, and the
constant $3/2$ of Lemma~\ref{contact-square}; the quantified loss is
discussed in Remark~\ref{rem-theta}.

\textbf{Verification protocol.} Every displayed identity, factorization,
completion of squares and ``$>0$'' assertion was re-checked in exact
arithmetic: rationals as exact fractions and $\sqrt3,\sqrt5$ in the
quadratic fields $\mathbb Q(\sqrt3)$, $\mathbb Q(\sqrt5)$, never in floating
point. Pointwise inequalities were additionally scanned on a
$2\times10^6+1$ point grid, and the finite-graph transfer was checked on all
$33864$ graphs of order $3\le n\le6$. The verification scripts are listed in
the accompanying audit report of 2 November 2026.

\begin{thebibliography}{99}
\raggedright
\bibitem{Weighted}
S. V. Borodachov, P. G. Boyvalenkov, P. D. Dragnev, D. P. Hardin,
E. B. Saff and M. M. Stoyanova,
Energy bounds for weighted spherical codes and designs via linear programming,
\emph{Analysis and Mathematical Physics} \textbf{15} (2025), article 19.
\href{https://doi.org/10.1007/s13324-024-01009-7}{doi:10.1007/s13324-024-01009-7}.
\bibitem{Hjoin}
D. M. Cardoso, H. Gomes and S. J. Pinheiro,
The $H$-join of arbitrary families of graphs,
\href{https://arxiv.org/abs/2101.08383v3}{arXiv:2101.08383v3} (2021).
\bibitem{DL}
B. Der\k{e}gowska and B. Lewandowska,
A simple proof of the Gr\"unbaum conjecture,
\emph{Journal of Functional Analysis} \textbf{285} (2023), 109950.
\href{https://doi.org/10.1016/j.jfa.2023.109950}{doi:10.1016/j.jfa.2023.109950}.
\bibitem{DL-SW}
B. Der\k{e}gowska and B. Lewandowska,
From Sidelnikov--Welch bounds to projection constants,
\href{https://arxiv.org/abs/2609.29422v1}{arXiv:2609.29422v1} (2026).
\bibitem{Ebrahimi}
J. Ebrahimi Boroojeni,
\emph{Sum of the Largest Eigenvalues of Symmetric Matrices and Graphs},
M.Sc. thesis, Simon Fraser University, 2008.
\bibitem{Glazyrin}
A. Glazyrin,
Moments of isotropic measures and optimal projective codes,
\href{https://arxiv.org/abs/1904.11159v2}{arXiv:1904.11159v2} (2020).
\bibitem{KLMPT}
H. Kumar, L. Liu, H. Monterde, S. Pragada and M. Tait,
Maximum spectral sum of graphs,
\href{https://arxiv.org/abs/2604.00512v2}{arXiv:2604.00512v2} (2026).
\bibitem{KMMP}
H. Kumar, B. Mohar, S. A. Mojallal and S. Pragada,
Stability of maximal relative projection constants,
\href{https://arxiv.org/abs/2609.03200v1}{arXiv:2609.03200v1} (2026).
\bibitem{LiLiu}
H. Li and X. Liu,
Sharp bounds for minimal dependencies of linear-form powers,
\href{https://arxiv.org/abs/2606.24349v1}{arXiv:2606.24349v1} (2026).
\bibitem{LS}
L. Lov\'asz and B. Szegedy,
Szemer\'edi's lemma for the analyst,
\emph{Geometric and Functional Analysis} \textbf{17} (2007), 252--270.
\bibitem{Mohar}
B. Mohar,
On the sum of $k$ largest eigenvalues of graphs and symmetric matrices,
\emph{Journal of Combinatorial Theory, Series B} \textbf{99} (2009), 306--313.
\href{https://doi.org/10.1016/j.jctb.2008.07.001}{doi:10.1016/j.jctb.2008.07.001}.
\bibitem{STW}
V. Sivashankar, Q. Tang and T. Wakhare,
Graph eigenvalues and projection constants,
\href{https://arxiv.org/abs/2608.02429v1}{arXiv:2608.02429v1} (2026).
\end{thebibliography}
\end{document}
